% 张三慧第四版V2第20章《电磁感应》，教材PDF物理页205--208。
% 习题20.1--20.27；图20.22--20.30均属原始条件，矢量重绘与逐图审核待办。
\section*{第20章\quad 电磁感应}
\subsection*{章末习题20.1--20.27}

\paragraph{习题 20.1}
图20.22：竖直无限长导线向上通 \(I=5\,\mathrm A\)，
右方有一水平短棒 \(ab\)，长 \(l=20\,\mathrm{cm}\)，
左端 \(a\) 距长导线 \(d=10\,\mathrm{cm}\)。
短棒平行长导线以速度
\(v=10\,\mathrm{m/s}\) 向上平移。
求棒两端动生电动势大小及哪端电势较高。

\paragraph{习题 20.2}
平均半径 \(R=12\,\mathrm{cm}\) 的
\(N=4\times10^3\) 匝线圈，
在 \(B=0.5\,\mathrm G\) 的地磁场内
以每秒 \(30\) 周的频率转动。
求线圈可产生的最大感应电动势，
以及达到该最大值的转轴和瞬时取向条件。

\paragraph{习题 20.3}
图20.23：无限长直线中通 \(I=5\,\mathrm A\)。
右侧共 \(N=10^3\) 匝的矩形线圈，
宽 \(a=10\,\mathrm{cm}\)、
沿直线方向高 \(L=20\,\mathrm{cm}\)，
以 \(v=2\,\mathrm{m/s}\) 远离长直线平移，
最近边距离 \(d=10\,\mathrm{cm}\)。
求该瞬间线圈的感应电动势。

\paragraph{习题 20.4}
沿用习题20.3的矩形线圈与长直线几何，
但线圈保持不动，
直导线电流按
\(i(t)=5\sin(100\pi t)\,\mathrm A\)
变化。求线圈的瞬时感应电动势。

\paragraph{习题 20.5}
图20.24：圆盘形区域半径 \(R\) 内有
垂直纸面向里的均匀变化磁场 \(B(t)\)；
盘内有长度 \(L\) 的静止金属直杆
\(ab\) 形成圆的一条弦，两端都在圆周上。
磁场随时间变化率 \(dB/dt\) 已知。
求杆两端之间沿杆的感生电动势及
两端哪一个电势高。
\textit{本题杆不运动，感生电场非保守，
不应直接套动生 \(B Lv\) 式。}

\paragraph{习题 20.6}
源图20.25为FBC-1“飞豹”飞机，
翼展 \(12.705\,\mathrm m\)，
在北京沿东西方向水平飞行。
北京地磁场竖直分量
\(B_z=0.42\times10^{-4}\,\mathrm T\)；
速度达 \(M=1.70\)（声速的1.70倍）。
求两翼尖之间动生电势差大小，
并按自东向西的飞行及地磁场竖直方向
判断哪端电势高。
\textit{声速需说明近似值。}

\paragraph{习题 20.7}
在海洋探测实验中，
海水穿过地磁场竖直分量
\(0.70\times10^{-4}\,\mathrm T\)。
在相距 \(200\,\mathrm m\) 的水流两端，
用高输入阻抗探针量得动生电势差
\(7.0\times10^{-3}\,\mathrm V\)。
求海水流速。

\paragraph{习题 20.8}
交流发电机由 \(N=120\) 匝矩形线圈组成，
尺寸 \(10.0\,\mathrm{cm}\times20.0\,\mathrm{cm}\)，
旋转轴处在圈平面内，
并垂直于外部均匀水平磁场
\(B=2.0\times10^{-2}\,\mathrm T\)。
要求输出感应电势幅值为
\(12.0\,\mathrm V\)，求每分钟转速。

\paragraph{习题 20.9}
一个 \(N\) 匝、面积 \(S\) 的小线圈
总电阻含灵敏电流计为 \(R\)，
将线圈从与待测磁场正交的初始取向
迅速移出磁场。
若流过的电荷总量为 \(q\)，
求待测 \(B\) 并解释为什么不依赖移出时间。

\paragraph{习题 20.10}
图20.26：金属圆盘厚 \(b\)、电阻率
\(\rho\)，角速度 \(\omega\) 绕中心轴转动，
在距轴约 \(r\) 的盘面上，
一个面积 \(a^2\) 的小正方形区域
有垂直盘面的磁感应强度 \(B\)。
推导旋转盘的阻力矩
和维持恒定转速的机械功率。
\textit{按小区域 \(a\ll r\)
及磁场几何近似使用局部涡流模型。}

\paragraph{习题 20.11}
电子感应加速器电子在半径 \(r\) 轨道上运行。
证明若电子轨道半径在加速中不变，
轨道处的磁感应强度 \(B_r\)
等于轨道内面积的平均磁感应强度
\(\overline B\) 的一半。
\textit{忽略辐射损失，需区分局部 \(B_r\)
与穿过轨道内部磁通的平均值。}

\paragraph{习题 20.12}
继续习题20.11，
在电子速率增加 \(\Delta v\) 很小时，
证明由变化磁场产生的切向电力
可与速度增加后额外需要的向心磁力
相匹配，使回旋轨道半径保持不变。
保留题定的一级近似假设。

\paragraph{习题 20.13}
长为 \(l\)、截面积
\(S=\pi R^2\) 的圆柱筒上
绕有两组均匀线圈，
匝数为 \(N_1,N_2\)。
忽略管外漏磁及端部效应，
求两线圈的互感系数。

\paragraph{习题 20.14}
一细圆线圈 \(a\) 有
\(N_a=50\) 匝、
截面积 \(S_a=4.0\,\mathrm{cm^2}\)，
同轴放在圆形线圈 \(b\) 中心，
\(N_b=100\) 匝、半径
\(R_b=20.0\,\mathrm{cm}\)。
求互感系数；
若线圈 \(a\) 中电流以
\(50\,\mathrm{A/s}\) 的速率减小，
求线圈 \(b\) 中磁通变化率和感应电动势。

\paragraph{习题 20.15}
空心长螺线管半径
\(r=2.0\,\mathrm{cm}\)，长
\(l=30.0\,\mathrm{cm}\)，
匀绕 \(N=1200\) 匝。
（1）求自感 \(L\)；
（2）电流按
\(dI/dt=3.0\times10^2\,\mathrm{A/s}\)
变化时，自感电动势大小。

\paragraph{习题 20.16}
长直螺线管通恒流
\(I=10.0\,\mathrm A\)，
通过每匝的磁通
\(\Phi=20\,\mu\mathrm{Wb}\)；
电流变化率
\(dI/dt=4.0\,\mathrm{A/s}\) 时
自感电动势大小
\(3.2\,\mathrm{mV}\)。
求螺线管的自感系数与总匝数。

\paragraph{习题 20.17}
图20.27：矩形截面、空气芯的环形螺线管，
截面高度 \(h\)、内半径 \(R_1\)、
外半径 \(R_2\)，共 \(N\) 匝。
（1）求自感；
（2）一根无限长直导线穿过环的中心轴，
求导线与环形螺线管间
\(M_{12}\) 和 \(M_{21}\)，证明二者相等。

\paragraph{习题 20.18}
两根平行长导线半径均为 \(a\)，
两导线轴线相距 \(D\)，
两电流大小相同、方向相反。
忽略导线内部磁场，
证明每单位长度的自感
\[
 \boxed{L'=\frac{\mu_0}{\pi}
 \ln\frac{D-a}{a}}.
\]

\paragraph{习题 20.19}
两个相互重叠、平面正交放置的
单匝平面线圈，忽略彼此互感，
各自自感为 \(L_1,L_2\)，
联接后的等效自感为 \(L\)。
证明串联时
\(L=L_1+L_2\)，
并联时
\(1/L=1/L_1+1/L_2\)。

\paragraph{习题 20.20}
图20.28：两个线圈自感
\(L_1,L_2\)，互感为 \(M\)。
（1）把2、3端相接、
以1、4端接电源，证明顺串联时
\(L=L_1+L_2+2M\)；
（2）把2、4端相接、
以1、3端接电源，证明反串联时
\(L=L_1+L_2-2M\)。

\paragraph{习题 20.21}
中子星表面磁场量级
\(10^8\,\mathrm T\)。
求磁场能量密度及按
\(E=mc^2\) 等效的质量密度，
以 \(\mathrm{kg/m^3}\) 表示。

\paragraph{习题 20.22}
实验室中普通磁场大小约
\(B=2.0\,\mathrm T\)，
强电场大小约
\(E=1.0\times10^6\,\mathrm{V/m}\)。
计算两种场的能量密度
\(u_B,u_E\)，并比较哪个有利于储能。

\paragraph{习题 20.23}
拟在超导线圈的强磁场中储存
\(1\,\mathrm{kW\,h}\) 能量。
若 \(B=1.0\,\mathrm T\)，
所需磁场体积多大？
若靠电流 \(I=500\,\mathrm A\)
储存该能量，线圈的自感系数需要多大？

\paragraph{习题 20.24}
一根长直铜导线横截面半径
\(r=5.5\,\mathrm{mm}\)，
通电 \(I=20\,\mathrm A\)，
电阻率
\(\rho=1.69\times10^{-8}\,\Omega\mathrm m\)。
求导线外贴近表面处电场和磁场
的能量密度并比较。

\paragraph{习题 20.25}
一根同轴电缆由中心实心金属圆柱
和外围薄圆筒组成，半径分别为
\(R_1,R_2\)，中间为空气或
\(\mu_r\simeq1\) 的介质。
中心芯导体通总电流 \(I\) 流出，
外筒通同样大电流返回。
假设芯横截面电流密度均匀，
求单位长度电缆储存的磁能，
并由 \(U_m=\frac12L I^2\)
求单位长度自感系数。

\paragraph{习题 20.26}
图20.29：两根半径 \(a\) 的
长直平行管状导体连接同一稳流电源，
中心间距 \(D_1\gg a\)，
电流一去一回。
忽略管内部磁场：
（1）求每单位长度电路自感；
（2）恒定电流 \(I\) 下固定一管，
另一管被磁斥力推至更大间距
\(D_2\)，求单位长度磁力所做功；
（3）求该过程单位长度磁场能量变化；
（4）判定感应电动势方向及电源的总功；
（5）解释磁力功、磁能变化、电源功之间的关系。

\paragraph{习题 20.27}
图20.30：两根很长的螺线管
截面积均为 \(S\)，一个同轴插入另一个。
每单位长度匝数为 \(n_1,n_2\)，
电流分别为 \(I_1,I_2\)。
证明两个螺线管的轴向相互作用力大小为
\[
 \boxed{F_m=\mu_0n_1n_2SI_1I_2}
\]
并说明同向电流产生吸引的条件。
